% Folder 57, batch 07 — archivists' pages 121–140.
% Pass: Opus 5.5 (claude-opus-5-5), 2026-10-03 — first pass, unchecked against the pages by a human.
%
% Skipped: 129, 131, 135 (typescript versos, cohomology drafts, nothing in his
% hand), 136 (blank). 124 and 140 are covers in his hand: their words are the
% section headings, without \page.
\documentclass[11pt,a4paper]{article}
\input{../preamble/grothendieck}

\folder{57}
\batch{7}
\pages{121}{140}
\foldertitle{Picard : tapuscrit annoté (s.d.), lettres (s.d., 1962).}
\dating{1962-[vers 1968]}
\watermark{Édition de démonstration}

\begin{document}

\section{Diviseurs à support dans les fibres}

\page{121}
\note{L'argument vient d'avant le lot : « la Proposition » invoquée ici, ses
fermés $X_i$, $Y_i$ et la fonction $f$ sont introduits dans les pages
précédentes. Page rapide, très raturée ; la prose de liaison est souvent
perdue.}

\textbf{Corollaire 3.} Sous les conditions de la proposition,
\struck{considérons} \add{soit} \struck{$D(X)$ un diviseur} $D$ un diviseur
\uncertain{positif} sur $X$, \struck{\ill{}} satisfaisant les conditions
suivantes :

(i) $D \cap X_f = 0$ ;

(ii) \struck{$D \cap \underline{D}_{X_i} = n(X_i)$} \add{$D \cap
\underline{D}_{Y_i} = n(Y_i)$}, i.e.\ le coefficient de $Y_i$ dans $D$
\uncertain{est indépendant de $i$}.

Alors $D = n\,\struck{\ill{}}\,\mathrm{div}(f)$.

\textbf{Corollaire 4} \struck{(Grandes conditions…)}. Sous les conditions de
la Prop., supposons \ill{} les conditions de \ill{} :

(ii) \uncertain{Pour tout} $X_\alpha$, $Y \cap X_\alpha$ est irréd.

(iii) $Y$ connexe \struck{\ill{}} « d'incidence » \ill{} : $Y_i$ et $Y_j$
\uncertain{sont reliés} \ill{} \emph{irréductibles} \uncertain{communes à}
$Y_i \cap Y_j$ \struck{\ill{} la \ill{}}.

Soit $D$ un diviseur \struck{\ill{}} sur $X$ tel que $D \cap X_f = 0$
\uncertain{en dehors de} $Y$ \ill{} \struck{$Y$ irréductible dans $X$}
\ill{}. Alors $D = n\,\mathrm{div}(f)$.

\marginal{N.B. les hypothèses \ill{} \uncertain{doivent} \ill{} : s'il existe
une \ill{} $Z_i$ \ill{} $Y_i \cap Y_j$ \ill{} $X \supset Z_i$ \ill{} \ill{} ;
alors \ill{} $Y = \ill{}$ \ill{}.}
\note{Note marginale de sa main, écrite en biais dans la marge gauche et
encadrée d'un trait ; on n'en lit que des fragments.}

\page{122}
Je \uncertain{dis} \ill{} \ill{} \uncertain{en} \uncertain{cas 3} \ill{} la
\uncertain{condition} \uncertain{pour} \struck{$X_i \cap X_j$} les \ill{}
\uncertain{soient} \uncertain{irréductibles} \add{\uncertain{irréductibles}}
\ill{} il \uncertain{suffit} \ill{} la \uncertain{condition}
\uncertain{précédente} (\ill{} la \uncertain{composante} $Y_i = Y \cap X_i$,
$Y_j = Y \cap X_j$) \struck{\ill{} $Y_i$}, $Y_j$ « irréductibles ».
\struck{$P$\ill{}}

Soit $X_{ij}$ une \uncertain{composante} irréductible de \struck{\ill{} le
\ill{}} $X_i \cap X_j$ \ill{} $Z_{ij}$ \uncertain{de} \ill{} \uncertain{dans}
$X_j$. \ill{} \ill{} $X = \operatorname{Spec} A$ ; \struck{\ill{}}
\uncertain{avec} $D$ \struck{\ill{}} de la \uncertain{forme}
$\frac{1}{n}\,\ill{}\,\mathrm{div}\,t$, $t \in$ \struck{$B$} $A(A)$.
\uncertain{Comme} $D \cap X_j = 0$, \ill{} $t \in A_j^{*}$, \struck{$x \in A$}
\struck{\ill{} $t = x f^{n}$, Soit $X_2(X_i)$, $f = f(f_i)$.}

\drawing{figure : les composantes $X_i$, $X_j$ découpées dans $X$ par des
arcs, $Y_i$ et $Y_j$ portés sur elles, $X_{ij}$ et $Z_{ij}$ marqués au point
d'intersection, avec flèches vers les étiquettes}

(iii) \uncertain{En} \ill{} \uncertain{restreignant} \ill{} $X_i$, \ill{}
\uncertain{ont} \uncertain{que}
\[
  t_i = \varepsilon_i f_i^{\,n_i}, \qquad \varepsilon_i \in A_i^{*},
\]
\struck{\ill{}} \ill{} \ill{}. \uncertain{Donc}, \uncertain{comme}
$Z_{ij} \cap X_j \neq \emptyset$, $t_i | Z_{ij}$ et $t_j | Z_{ij}$
\ill{}.

Soit $t_i = t|X_i$, $f_i = f|X_i$ \uncertain{\ill{}}
\[
  t_i = \varepsilon_i f_i^{\,n_i} \qquad (\varepsilon_i \in A_i^{*},\
  n_i \in \mathbb{Z}).
\]
\struck{\ill{}} $Z_{ij} \subset X_{ij}$ \struck{$X_{ij}$} \ill{},
$t_i|X_{ij}$ et $t_j|X_{ij}$ constants, \ill{} soit
$t_{ij} = t_i|\ill{}$ \ill{} \struck{\ill{}} \ill{} \ill{} \struck{$t_{ij} =
\varepsilon_{ij} f$} $f_{ij} = f_{i}|\ill{}$ \struck{\ill{}},
$\varepsilon_{ij} = \varepsilon_i|Z_{ij}$, $\varepsilon_{ji} = \varepsilon_j|Z_{ij}$.
\ill{} \ill{} \uncertain{restrictions} \uncertain{sur} $Z_{ij}$ \ill{}

\marginal{\uncertain{Attention}, \ill{} \uncertain{n'est pas} \ill{}
$X_i \times_X Y$ \uncertain{réduit} ; \ill{} $V(\ill{})$ \ill{}. — \ill{}
$X_{ij}$.}

\page{123}
On aura \uncertain{évidemment} \ill{}
\[
  t_{ij} = \varepsilon_{ij} f_{ij}^{\,n_i} = \varepsilon_{ji} f_{ji}^{\,n_j},
\]
\struck{$\varepsilon_{ij} f^{n_i} = \varepsilon_{ji} f^{n_j}$} \struck{Supp.
Je dis} \struck{\ill{}} Je dis \uncertain{que} $n_i = n_j$. \uncertain{Sinon},
\ill{} \ill{}, \uncertain{si} $n_i > n_j$, \ill{} \uncertain{on aurait}
\struck{$(\varepsilon_{ij}/\varepsilon_{ji}) f_{ij}^{\,n_i - n_j} = 1$}
\struck{et \ill{}}
\[
  f_{ij}^{\,(n_i - n_j)} = \varepsilon_{ji}/\varepsilon_{ij},
\]
donc $f_{ij}$ \uncertain{serait} \uncertain{une} unité, \ill{}
\uncertain{puisque} \struck{\ill{} pas \ill{}} \ill{} \uncertain{sur}
$Z_{ij}$, \uncertain{puisque} \ill{} $f_{ij}$ s'\ill{} \ill{}
\uncertain{sur} $Z_{ij} \subset X_{ij}$.

\textbf{Théorème.} Soit $f \colon X \to Y$ un \struck{\ill{}} morphisme
propre \add{et} \emph{\uncertain{séparable}},
\uncertain{$Y$ loc.\ noeth.}\add{\uncertain{tel que} \ill{} les \ill{}
$f^{-1}(y)$} \struck{\ill{}} \ill{} \ill{} \uncertain{composantes}
irréductibles \emph{\uncertain{géométriques}} de $f^{-1}(y)$, et
\uncertain{\ill{}} $n(y)$ \uncertain{loc.} \ill{} \uncertain{sur} $X$.
\struck{Alors} (i) Soit $\underline{L}$ un faisceau inversible.
(ii) Il existe un \uncertain{plus grand} \struck{\ill{}} \uncertain{sous-préschéma}
\struck{\ill{}} \uncertain{fermé} $Y'$ de $Y$ tel que
$\underline{L} \otimes_Y Y'$ \uncertain{soit} \uncertain{isomorphe} \ill{}
\uncertain{à un faisceau} \uncertain{venant} \struck{\ill{}} \ill{}
\uncertain{sur} $Y$, \uncertain{c'est} \ill{} \struck{\ill{}}
$\mathcal{S}(\underline{L}/X/Y)$.

(ii) \ill{}.
\[
  \mathcal{S}(\underline{L}_1, X_1, Y_1) \simeq
  \mathcal{S}(\underline{L}/X/Y) \times_Y Y_1,
\]
\note{La page s'arrête sur cette formule ; la suite du théorème n'est pas
dans le lot.}
\note{Le sigle $\mathcal{S}(\underline{L}/X/Y)$ est lu sur une lettre
cursive bouclée ; la lecture de la lettre est incertaine.}

\section{Th.\ d'existence relatifs}
\note{Intitulé de sa main, seul sur la page 124, qui sert de chemise à ce qui
suit.}

\page{125}
\marginal{($p \colon X' \to X$ morphisme \ill{}. (\ill{} \uncertain{pour}
$X'' = X' \times_X X'$, $X''' = X' \times_X X' \times_X X'$, et \ill{}
\uncertain{désigne} \uncertain{par} $f''$, $f'''$ \uncertain{les morphismes}
\ill{}).}

\textbf{Proposition.} Soient $f \colon X \to S$ et $f' \colon X' \to S'$
\uncertain{deux} morphismes \ill{} \uncertain{supposés} \ill{}
$f = (X \xrightarrow{f_1} \Sigma \xrightarrow{f_2} S)$, avec $f_2$ \ill{}
\struck{et $f_1(\mathcal{O}_X) \simeq \mathcal{O}_\Sigma$}, et de même
$f' = (X' \xrightarrow{f'_1} \Sigma' \xrightarrow{f'_2} S)$. \uncertain{On
suppose} donnée une section $g$ de $X$ sur $\Sigma$, une section $g'$ de $X'$
sur $\Sigma'$, et \uncertain{on} désigne par $Z$ le \uncertain{sous-préschéma}
(fermé) $g(\Sigma)$ de $X$, par $Z'$, $Z''$ \ill{} \ill{}, \add{$Z''$}
\uncertain{les images inverses} \uncertain{dans} $X'$, $X''$ \add{de
$Z'' = Z' \times_Z Z'$} \struck{\ill{}, \uncertain{posons} (i)
$g(\Sigma') = Z'$. \ill{} $X'' = X' \times_X X'$, $X''' = X' \times_X X'
\times_X X'$, et \ill{} \uncertain{désigne} par} \add{$f''$, $f'''$ :
$\Sigma' \to S$, $Z'' \to S$}.

(i) \struck{\ill{}} $f_1$, \struck{$\Sigma$} $f'_1$, \struck{$\Sigma$}
\struck{\ill{}}, $Z'' \to S$ \add{sont plats}.
\struck{[\ill{} $f, f'$, $Z'' \to S$, $Z''' \to S$]} $f_1$ et $p$ propres.

(ii) $f_{1*}(\mathcal{O}_X) \simeq \mathcal{O}_\Sigma$ et
$f'_{1*}(\mathcal{O}_{X'}) \simeq \mathcal{O}_{\Sigma'}$.

\struck{\emph{universellement} \ill{} $S$.}

(iii) $Z' \to Z$ est un morphisme de descente
\uncertain{universel} pour les \uncertain{Modules} inversibles.

(iv) $p \colon X' \to X$ est un morphisme de descente \emph{effective}
\uncertain{universel} pour les Modules inversibles.

(v) \struck{$g$} $g' \colon \Sigma' \to X'$ \uncertain{est une} \uncertain{section}
\ill{} \uncertain{sur} $Z'$.

Sous ces conditions, $\underline{\mathrm{Pic}}_{X/S} \to
\prod_{\Sigma/S} \underline{\mathrm{Pic}}_{X/\Sigma}$ \ill{} \struck{\ill{}}
$\underline{\mathrm{Pic}}_{X'/\Sigma'}$ \uncertain{est} représentable par
des morphismes \emph{affines}.

\textbf{Cor.} Si $\underline{\mathrm{Pic}}_{X'/\Sigma'}$ \struck{est
représentable} \add{existe}, il en \uncertain{est} de même de
$\underline{\mathrm{Pic}}_{X/S}$, et $\underline{\mathrm{Pic}}_{X/S} \to
\underline{\mathrm{Pic}}_{X'/S'}$ est un morphisme affine.

\page{126}
\struck{Comme} $f_1$ \uncertain{est} propre et
$\mathcal{O}_\Sigma \simeq f_*(\mathcal{O}_X)$ \uncertain{universellement},
il s'ensuit \struck{\ill{}} que $\underline{\mathrm{Pic}}_{X/S}(T)$
s'interprète par les $\underline{L}$ inversibles sur $X_T$
$g_T$-rigidifiés. \uncertain{Analogue} pour
$\underline{\mathrm{Pic}}_{X'/S'}(T)$. \uncertain{Utilisant} \ill{}
\uncertain{la} \uncertain{théorie} \uncertain{de la} \uncertain{descente} \ill{}, \ill{} \ill{}
\marginal{$f_1$, $f'_1$ plats (i) ; $\Sigma \to S$, $\Sigma' \to S$ plats ;
(iii).}

\marginal{L'homomorphisme $\underline{\mathrm{Pic}}_{X/S} \to
\underline{\mathrm{Pic}}_{X'/S'}$ s'interprète \uncertain{alors} \uncertain{de
façon} \uncertain{simple} \uncertain{grâce} \uncertain{aux} $g' \colon
\Sigma' \to X'$ \uncertain{qui} \struck{\uncertain{majorent}} \uncertain{sur}
$Z'$ (v).}

\uncertain{revient} à \uncertain{ceci} : \emph{Soit} $\underline{L}'$ sur
$X'$ \emph{avec} une \emph{$g'$-rigidification} $\psi'$
\struck{\uncertain{Représentable}} \emph{\uncertain{Prouvons que}}
\emph{le foncteur} $T \mapsto$ \uncertain{ens.} des $\underline{L}$
$g$-rigidifiés sur $X_T$, \uncertain{tels que} $p^{*}(\underline{L}) \simeq
\underline{L}'$ \emph{\uncertain{comme} $g$-rigidifiés}, \emph{est
représentable par un schéma affine sur $S$}.

\uncertain{Or} la donnée de $\underline{L}$ \emph{\uncertain{équivaut à}} :

1°) Un isom.\ $q_{1T}^{*}(\underline{L}'_T) \xrightarrow{\ \varphi\ }
q_{2T}^{*}(\underline{L}'_T)$, où $q_{1T}^{*}(\underline{L}'_T) =
q_1^{*}(\underline{L}')_T$ et $q_{2T}^{*}(\underline{L}'_T) =
q_2^{*}(\underline{L}')_T$ ;

2°) Un isom.\ $\underline{O}_{Z'} \xrightarrow{\ \psi\ }
\struck{\ill{}}\, i'^{*}(\underline{L}'_T)$, où $i'^{*}(\underline{L}'_T) =
i'^{*}(\underline{L}')_T$ ;

satisfaisant les conditions suivantes :

a) Transitivité : $r_{2T}^{*}(\varphi)\, r_{1T}^{*}(\varphi) =
r_{3T}^{*}(\varphi)$ ;

b) $\bar{q}_{2T}^{\,*}(\psi) = i''^{*}_{T}(\varphi)\,
\bar{q}_{1T}^{\,*}(\psi)$ ;

c) \uncertain{Condition} \ill{} \uncertain{sur}
$\bar{g}'^{\,*}_{T}$\ill{}.
\note{La condition c) est coupée au bord droit de la feuille ; on n'en lit
que le début.}

\struck{On \ill{} de}

Le diagramme de ses notations, en bas à gauche de la page :

\begin{tikzcd}[column sep=small, row sep=small, nodes={font=\scriptsize}]
X & X' \arrow[l, "p"'] & X'' \arrow[l, "{q_1, q_2}"'] & X''' \arrow[l, "{r_1, r_2, r_3}"'] \\
Z \arrow[u, "i"] & Z' \arrow[u, "{i'}"] \arrow[l, "\bar{p}"] & Z'' \arrow[u, "{i''}"] \arrow[l, "{\bar{q}_1, \bar{q}_2}"'] & \\
\Sigma \arrow[u, "g"] & \Sigma' \arrow[u, "{\bar{g}'}"'] & &
\end{tikzcd}

\note{Diagramme recomposé : $X''$ est relié à $X'$ par deux flèches $q_1,
q_2$, $X'''$ à $X''$ par trois flèches $r_1, r_2, r_3$, que je réduis chaque fois à une
flèche étiquetée ; de même pour $\bar{q}_1, \bar{q}_2$ de $Z''$ vers $Z'$ ; au-dessus, trois arcs étiquetés $s_1, s_2, s_3$ reviennent
de $X'''$ vers $X'$ ; un astérisque surmonte $X''$. Les étiquettes des
flèches verticales de la dernière ligne sont en partie raturées.}

\page{127}
Or les données 1° et 2° s'expriment \uncertain{par} un
\uncertain{homomorphisme} de $T$ dans un $S$-préschéma \struck{affine}
$K$, $L$ sur $S$ affine sur $S$. Les \struck{morphismes} \ill{}
\uncertain{isomorphismes} de $s_1^{*}(\underline{L}')_T$ dans
$s_2^{*}(\underline{L}')_T$, \ill{} de \struck{$(q_1 i'')^{*}(L')$}
$(\ill{})^{*}_{T}$ dans \ill{} \ill{} de $\underline{O}_{\Sigma}$ dans
$(i'\bar{g}')^{*}(\underline{L}')$, \uncertain{et} \ill{}
$(q_2 i'')(\underline{L}')_T$, \struck{s'expriment également \ill{} par des
morphismes} \ill{} \ill{} \ill{} \uncertain{par} \ill{} \uncertain{S-schémas}
$M$, $N$ \ill{} affines \struck{sur $S$, et \ill{} \ill{} \ill{} de
$(i'\bar{g}')^{*}(\underline{L}')$ sur $\Sigma'$} s'expriment par les
$S$-morphismes \ill{} \ill{} \ill{} $S$-préschémas $M$, $N$, $P$
\struck{\ill{}} affines sur $S$, \uncertain{au-dessus} \uncertain{de} $S$.

\marginal{$f''$ et $Z' \to S$ plats (i) ; $\Sigma' \to S$, $f'''$ et
$Z'' \to S$ plats (i).}

Les conditions écrites a), b), c) s'expriment par le fait que \struck{deux}
\add{couples} \uncertain{de} morphismes
\[
  K \rightrightarrows M, \qquad L \rightrightarrows N, \qquad
  L \rightrightarrows P
\]
s'égalisent. \uncertain{Donc} \uncertain{si} \struck{$K$, $L$} $U$, $V$
\ill{} \uncertain{les} \uncertain{noyaux} \uncertain{des} \ill{}
$K \rightrightarrows M$ \uncertain{et} $L \rightrightarrows N \times P$
\uncertain{resp.} [\uncertain{qui sont} \uncertain{finis} \uncertain{sur}
$K$, $L$, \uncertain{donc} affines sur $S$], le foncteur
\uncertain{envisagé} \uncertain{est} représenté par $U \times V$, cqfd.

\page{128}
\note{Page barrée de trois traits obliques, mais lisible ; la rature vaut
pour le premier paragraphe, le corollaire qui suit est hors des traits.}
\uncertain{Or} les données 1°, 2° \uncertain{correspondent}
\uncertain{respectivement} \uncertain{à des} \struck{\uncertain{sections}}
\add{\uncertain{sections}} de préschémas affines
\add{[\uncertain{puisque} \uncertain{on suppose} \uncertain{que} $Z'$,
$X''$ \add{propres et} plats sur $S$]}
\struck{$K$, $L_T$} sur $S$, \uncertain{dans} 1°, 2° : \uncertain{une}
section \uncertain{de} $\struck{K} \times \struck{L}$ \uncertain{sur} $S$.
\struck{\ill{}} D'autre part, les isom.\ de $(L''_1)_T$ sur $(L''_3)_T$
\uncertain{correspondant} \uncertain{aux} \struck{sections d'un $\ill{}$ sur}
\ill{} $M_T$ affine, \struck{sur $S$} \ill{} \uncertain{à ceux} de
$\underline{O}_{Z''_T}$ \uncertain{dans} $L''^{*}_{T}(L''_1)$ \uncertain{aux}
sections d'un $N_T$ affine [\uncertain{car} $X'''$, $Z''$ \struck{\ill{}}
\uncertain{propres} et plats sur $S$] \struck{\ill{}}. Les conditions a) et b)
s'écrivent donc \uncertain{en} \struck{disant} \uncertain{que} \ill{}
\uncertain{coïncident} \uncertain{des} morphismes donnés
\[
  K \rightrightarrows M, \qquad L \rightrightarrows N, \qquad
  \text{i.e.} \quad K \times M \rightrightarrows L \times N .
\]
\uncertain{Donc} \ill{} \uncertain{donné} \uncertain{par} un
\uncertain{sous-schéma} \uncertain{fermé} de $K \times M$. \uncertain{Cela}
\uncertain{achève} la démonstration.

\textbf{Corollaire.} $X$ séparable et complet sur le corps $k$.
\struck{Supposons les $\underline{\mathrm{Pic}}_{X_i/k}$ existent.}
\struck{Alors $\underline{\mathrm{Pic}}_{X/k}$ existe.} Soient $X_i$
\uncertain{les} \uncertain{composantes} \add{\uncertain{(réduites)}}
\struck{\ill{} des $X_i$ irréd.\ de $X$, \ill{}} \uncertain{irréductibles},
\uncertain{alors} \uncertain{le} \struck{\ill{}}
\[
  \underline{\mathrm{Pic}}_{X/k} \longrightarrow
  \prod_i \underline{\mathrm{Pic}}_{X_i/k}
\]
\struck{\ill{} le noyau \ill{} des \ill{} morphismes} \struck{est affine}
\uncertain{est} représentable par des morphismes affines. \uncertain{Donc}
\uncertain{si} les $\underline{\mathrm{Pic}}_{X_i/k}$ existent, \uncertain{il
en est de même de} $\underline{\mathrm{Pic}}_{X/k}$.

\page{130}
\uncertain{Résultats} \uncertain{plus} \uncertain{élémentaires}
\struck{\ill{}} \ill{} \ill{}

\textbf{Lemme.} $X$ \struck{séparable} \add{\uncertain{séparable}} propre sur
le corps $k$, et \struck{soit} $X_1$ et $X_2$ deux parties fermées de $X$,
séparables sur $k$, telles que $X = X_1 \cup X_2$. Alors
\[
  \underline{\mathrm{Pic}}_{X/k} \longrightarrow
  \underline{\mathrm{Pic}}_{X_1/k} \times \underline{\mathrm{Pic}}_{X_2/k}
\]
est représentable par \uncertain{un} morphisme affine.

\struck{Si $X_1 \cap X_2 = \emptyset$, \ill{} \ill{} \ill{} \ill{}. Supposons
$X_1 \cap X_2 \neq \emptyset$.}

\uncertain{Comme} les \uncertain{foncteurs} \ill{} \uncertain{sont}
\uncertain{compatibles} \uncertain{avec} \struck{\uncertain{passage}} \uncertain{les}
\uncertain{descentes} \uncertain{fid.\ plates} \uncertain{quasi-compactes},
et que \struck{les} \uncertain{morphismes} \ill{} \uncertain{affines}
\uncertain{sont} \ill{} \uncertain{de} \uncertain{morphismes} \ill{}
\uncertain{descente} \uncertain{effective} \uncertain{des} morphismes affines,
\uncertain{rien} \uncertain{n'est} \uncertain{changé} \uncertain{pour} \ill{}
\uncertain{en} \uncertain{faisant} \uncertain{un} \uncertain{changement}
\uncertain{de base} $k'/k$, \ill{} \uncertain{finissant} \uncertain{par} \ill{}
\uncertain{que} $k'/k$ est une extension de corps. \uncertain{Donc}
\uncertain{on peut} \ill{} \ill{} \uncertain{que} \ill{}
\uncertain{dans} \uncertain{ce} \uncertain{cas} \ill{} \uncertain{on a}
\struck{$F$} $= X_1 \cap X_2$ \ill{} \uncertain{a} \uncertain{un} \ill{}
\uncertain{rat.} \uncertain{sur} $k$, i.e.\ une section \ill{} \uncertain{sur}
$\operatorname{Spec}(k)$. \uncertain{On} \uncertain{applique} \uncertain{alors}
le \uncertain{théorème} \uncertain{avec} $X' = X_1 \amalg X_2$.

\struck{\ill{}}

\textbf{\uncertain{Attention}} \uncertain{comme} \ill{} \uncertain{dans} le
\uncertain{corollaire} : \uncertain{on se} \uncertain{ramène} \uncertain{au}
\uncertain{cas} où $k$ \uncertain{est} \uncertain{alg.\ clos}, et $X$
\uncertain{connexe}, \ill{} \uncertain{alors} \ill{} \emph{\ill{}} \ill{}
$X_1$ et $X_2$ \uncertain{sont} \emph{\uncertain{connexes}} et
$X_1 \cap X_2$ \uncertain{a} \uncertain{un} $p^t$ \uncertain{rat.}\ $/k$.
\uncertain{Pour} \ill{} \uncertain{dans} \ill{} \uncertain{les} \ill{}
\note{La dernière ligne, au bord inférieur roussi de la feuille, se perd ;
l'argument continue hors du lot.}

\marginal{$\underline{\mathrm{Pic}}_X \to \prod \underline{\mathrm{Pic}}_{X_i}$
\ill{} \uncertain{affine} $\to \underline{\mathrm{Pic}}_{X'} \times
\underline{\mathrm{Pic}}_{X''}$ \uncertain{dans} \uncertain{le} \uncertain{cas}
\ill{}.}

\page{132}
\textbf{Théorème.} Soit \struck{$X' \to X''$} \add{$f \colon X \to Y$} un
morphisme \emph{surjectif} de préschémas \uncertain{propres} \uncertain{sur}
\uncertain{un} corps $k$. Alors le noyau $\underline{N}$ de
$\underline{\mathrm{Pic}}_{Y/k} \to \underline{\mathrm{Pic}}_{X/k}$ est
\struck{contenu} \uncertain{extension} \uncertain{d'un} \uncertain{groupe}
\uncertain{fini} \uncertain{par} un groupe \emph{affine} $G$ \uncertain{sur} $k$.
\note{Un point d'interrogation, de sa main, dans la marge gauche face à
l'énoncé.}

\textbf{Corollaire} \add{\uncertain{Donc}} \uncertain{si}
$\underline{\mathrm{Pic}}_{Y/k}$ et $\underline{\mathrm{Pic}}_{X/k}$ existent,
le noyau de $\underline{\mathrm{Pic}}_{Y/k} \to \underline{\mathrm{Pic}}_{X/k}$
est un groupe \uncertain{quasi-affine}.

\textbf{Corollaire.} Sous les conditions précédentes, \ill{}
$\mathrm{NeuSev}_0(Y/k) \to \mathrm{NeuSev}_0(X/k)$ est
\emph{injectif}, et si $A$ et $B$ sont les « \uncertain{parties}
\uncertain{abéliennes} » de $\underline{\mathrm{Pic}}_{Y/k}$ et
$\underline{\mathrm{Pic}}_{X/k}$ [\uncertain{définies} \uncertain{à}
\uncertain{isogénie} \uncertain{près}], alors $A \to B$ a un noyau fini.
\note{« NeuSev » : abréviation lue telle quelle (Néron--Severi) ; le sous-indice
$_0$ est incertain.}

\uncertain{\textbf{Dém.}} \struck{$\underline{\mathrm{Pic}}_X \to
\underline{\mathrm{Pic}}$}, on \uncertain{peut} \uncertain{pour} \uncertain{simplifier}
\uncertain{supposer} \uncertain{que} \uncertain{tous} les
$\underline{\mathrm{Pic}}$ \uncertain{envisagés} \uncertain{sont}
représentables. a) On \uncertain{peut} supposer $k$ \emph{alg.\ clos}
(\uncertain{ou} \uncertain{non} ?). b) On a un diagramme de schémas
\emph{\uncertain{commutatif}}

\begin{tikzcd}[column sep=small, row sep=small]
\underline{\mathrm{Pic}}_{X/k} \arrow[d, no head] & \underline{\mathrm{Pic}}_{Y/k} \arrow[l] \arrow[d, "\text{affine}"] \\
\underline{\mathrm{Pic}}_{X_{\mathrm{red}}} & \underline{\mathrm{Pic}}_{Y_{\mathrm{red}}} \arrow[l, "?"']
\end{tikzcd}

\note{Le carré est fermé à droite par une accolade suivie de
« \uncertain{dont} » ; la page s'arrête là. Le trait vertical de gauche est
tracé sans pointe ; l'étiquette « affine » de la flèche verticale de droite est d'une lecture incertaine.}

\page{133}
\note{Feuillet dactylographié (un exposé sur le complexe de Koszul et la
cohomologie de Čech d'un recouvrement ouvert, § 2) qui n'est pas transcrit ;
seule l'annotation marginale à la main, face à la Proposition 1 et à son
corollaire, l'est.}
\marginal{\uncertain{Ajouter un corollaire} : les $f_i$ \uncertain{annulent}
les $H_{\cdot}(\underline{f}, M)$, $H^{\cdot}(\underline{f}, M)$.}

\page{134}
\struck{\uncertain{Il convient} \ill{} \ill{} \uncertain{factorisation} \ill{}
\ill{}.}

\[
  X \to Y' \to Y
\]
\note{Cette factorisation est encadrée et rattachée par un trait à la ligne
biffée au-dessus.}

\uncertain{Cela} \uncertain{nous} \uncertain{ramène} \uncertain{au} \uncertain{cas}
\uncertain{où} $Y$ \uncertain{est} \emph{réduit}.

c) \uncertain{Prenons} la factorisation de Stein
\[
  X \to Y' \to Y ;
\]
\uncertain{comme} $Y' \to Y$ \uncertain{est} \uncertain{fini} surjectif,
\uncertain{donc} un \uncertain{épimorphisme} \uncertain{fini}
\uncertain{factorisable} \uncertain{en} \uncertain{produit}
d'\uncertain{épimorphismes} \uncertain{stricts}. \uncertain{Par}
« \uncertain{Seshadri} », $\underline{\mathrm{Pic}}_Y \to
\underline{\mathrm{Pic}}_{Y'}$ \uncertain{est} \uncertain{affine}. \uncertain{On}
\uncertain{est} \uncertain{ramené} \struck{\uncertain{à prouver} que
$\underline{\mathrm{Pic}}_{Y'} \to \underline{\mathrm{Pic}}_X$ \ill{},
\ill{}} \ill{} \uncertain{au} \uncertain{cas} \uncertain{où}
$f_*(\mathcal{O}_X) \simeq \mathcal{O}_Y$.

d) \uncertain{Mais} \uncertain{dans} \uncertain{ce} \uncertain{cas} \ill{}
\ill{} \ill{} $\underline{\mathrm{Pic}}_X \to \underline{\mathrm{Pic}}_Y$
\uncertain{est} \ill{} \ill{}, \ill{} \ill{} \uncertain{fermeture} \uncertain{dans}
\uncertain{affine}.

\marginal{\uncertain{Attention} \uncertain{au} \uncertain{fait} \uncertain{que}
$\underline{\mathrm{Pic}}_X$ \add{et $\underline{\mathrm{Pic}}_Y$}
\struck{\ill{}} \ill{} \ill{} \ill{} \ill{} \ill{} \ill{}, \ill{} \ill{}
\uncertain{affine} \uncertain{réd.}!}

N.B. \uncertain{On} \uncertain{s'est} \uncertain{servi} \uncertain{de}
l'existence \uncertain{de} $\underline{\mathrm{Pic}}_Z$ \uncertain{pour}
$Z = Y_{\mathrm{red}}$, $X_{\mathrm{red}}$, et \uncertain{pour} \uncertain{les}
préschémas finis \uncertain{sur} $Y_{\mathrm{red}}$ ; \struck{\ill{}}
\uncertain{on} \uncertain{peut} \uncertain{s'en} \uncertain{passer}
\uncertain{sans} \uncertain{doute} \uncertain{(?)} \ill{}].
[\uncertain{cf.} $Z \to G$, $G = \mathrm{Grass}$ \uncertain{ou} \uncertain{VA}
\ill{}]

\page{137}
Soit $f \colon X \to S$ un morphisme se factorisant en
\[
  X \xrightarrow{\ \varphi\ } \Sigma \xrightarrow{\ \psi\ } S,
\]
avec $\psi$ fini, et $\varphi$ tel que
$\underline{O}_\Sigma \simeq \varphi_*(\mathcal{O}_X)$.

1) \uncertain{Alors} \uncertain{on a}
\[
  0 \to \mathrm{Pic}(\Sigma) \to \mathrm{Pic}(X) \to \mathrm{Pic}'(X/S)
\]
[suite exacte provenant de la suite spectrale de Leray].
\struck{\ill{}} \emph{Si} $X/\Sigma$ admet une section, le dernier
homom.\ est un \uncertain{épim.}, i.e.\ dans ce cas
\[
  \mathrm{Pic}'(X/S) \simeq \mathrm{Pic}(X)/\mathrm{Pic}(\Sigma).
\]
\note{Dans la formule, $X$ et $\Sigma$ sont récrits par-dessus eux-mêmes.}

Ce dernier fait se voit ainsi : \struck{Pour} Soit $g \colon \Sigma \to X$
une section de $X/\Sigma$ ; alors

(i) \struck{\ill{}} \uncertain{Les} Modules inversibles $g$-rigidifiés sur
$X$ \uncertain{n'ont} \uncertain{pas} \uncertain{d'automorphismes}
\uncertain{non} triviaux, (ii) \uncertain{Tout} Module inversible
$\underline{L}$ sur $X$ \uncertain{est} \struck{\ill{}}
$\mathrm{Pic}'(X/S)$-\emph{\uncertain{équivalent}} \uncertain{à}
\struck{l'image} un Module inversible $g$-rigidifié
[\uncertain{savoir} $\underline{L} \otimes \varphi^{*}(g^{*}(\underline{L}))$]
\note{sic, sans l'inverse qu'on attendrait sur $g^{*}(\underline{L})$}
[\uncertain{définissant} : isom.\ \uncertain{unique} \uncertain{près}].
[\uncertain{On} \uncertain{utilise} le fait que $\mathrm{Pic}'(\Sigma/S) = 0$].

\uncertain{On} \uncertain{le} \uncertain{prouve} \uncertain{par} la suite
\uncertain{spectrale} \uncertain{de} \uncertain{Leray} \uncertain{dir.}\ \ill{}

\page{138}
\note{Feuille écrite tête-bêche ; la partie lisible est barrée par deux
grands arcs et un trait horizontal, et abandonnée.}
\struck{Supposons (i) $Z$ fini sur $S$, (ii) $f_*(\mathcal{O}_X)
\xrightarrow{\sim} \varphi_*(\underline{O}_S)$.}
\struck{Alors (i) Tout Module inversible sur $X$ est}
\struck{$\mathrm{Pic}'(X/S)$-équivalent à un Module}

\drawing{à gauche, deux petits diagrammes : $Z \xrightarrow{i} X$, $Z$
fermé, avec $f \colon X \to S$ et $\varphi \colon Z \to S$ ; et le triangle
$X \to \Sigma \to S$, $X \to S$ ; un troisième croquis du même triangle, isolé
au bas de la page}

\page{139}
2) Soit $S' \to S$ \uncertain{fid.\ plat} et \uncertain{quasi-compact}. Alors

(i) $\mathrm{Pic}'(X/S) \to \mathrm{Pic}'(X'/S')$ \uncertain{est}
injectif. (ii) Si $X/\Sigma$ \uncertain{a} \struck{\ill{}} une section
\struck{\ill{}} localement sur $S$, alors
\[
  \mathrm{Pic}'(X/S) \to \mathrm{Pic}'(X'/S') \rightrightarrows
  \mathrm{Pic}'(X''/S'')
\]
est exact.
\marginal{\uncertain{généralisées} \uncertain{à} \uncertain{des}
\uncertain{notions} \ill{} \uncertain{de} \uncertain{l'exposé} \ill{}}

\uncertain{La} \uncertain{preuve} \uncertain{de} (i) \uncertain{\ill{}}
\uncertain{d'un} $\underline{L}$ \uncertain{sur} $X$. \add{\uncertain{En
vertu de} 1)} \uncertain{L'hyp.}\ \uncertain{signifie} \uncertain{que}
$\underline{L}'$ \uncertain{provient} \uncertain{de} $\Sigma'$. \uncertain{Or}
$\Sigma'$ \uncertain{est} \uncertain{fid.\ plat} \uncertain{quasi-compact}
\uncertain{sur} $\Sigma$. \uncertain{Cela} \uncertain{fait} \uncertain{que}
$\underline{L}$ \uncertain{provient} \uncertain{de} $\Sigma$.

\uncertain{Pour} (ii), \add{\uncertain{qui est}} \uncertain{la}
\uncertain{question} \uncertain{est} \uncertain{locale} \uncertain{sur} $S$,
\uncertain{on} \uncertain{peut} \uncertain{donc} \uncertain{supposer} \uncertain{que}
$X$ \uncertain{a} \uncertain{une} section \uncertain{sur} $\Sigma$. \uncertain{Alors}
\uncertain{on} \uncertain{utilise} la \uncertain{théorie} \uncertain{de}
\uncertain{descente} \uncertain{comme} \uncertain{dans} \uncertain{le}
raisonnement \uncertain{connu}.

\textbf{Corollaire.} \uncertain{Sous} les conditions précédentes,
\[
  \mathrm{Pic}'(X/S) \to \mathrm{Pic}^{*}(X/S)
\]
\uncertain{est} injectif, et \uncertain{est} bijectif \uncertain{si} $X/\Sigma$
\struck{\ill{}} \uncertain{a} une section loc.\ \uncertain{sur} $S$.
\uncertain{Si} $X/\Sigma$ \uncertain{a} une section, \uncertain{on a} un isom.
\[
  \mathrm{Pic}(X/S) \simeq \mathrm{Pic}(X)/\mathrm{Pic}(\Sigma).
\]

\section{Picard des schémas en courbes}
\note{Intitulé de sa main, seul sur la page 140, qui sert de chemise à la
suite du dossier, hors du lot.}

\end{document}
